\chapter{The orchids (Orchidaceae)}

In this family the germination of seeds and the growing on of the seedlings is a
highly specialized field. Propagation of hardy temperate zone orchids has been
achieved in the past few years, both by meristem methods and from seeds. The
following summary is based on presentations at a symposium sponsored by the
Brandywine Conservancy and held March 11-12 1989, at the Brandywine Museum in
Chadds Ford, Pennsylvania.

First a plea to refrain from purchasing collected native orchids and to refrain
from personally collecting them unless one is willing to make the requisite
committment to studying their culture. Only a few species (Cypripedium andrewsii,
calceolus, and candidum for example) will grow in neutral garden loams and even
these are susceptible to fatal viral and fungal diseases. Paul Keisling in Mass, has
had the most success, and he uses special microclimates and plenty of fungicide. At
present there are only two commercial establishments in the U.S. that purport to
propagate their own orchids and both operations are more in the nature of cropping
natural populations on their own lands. It would be possible to propagate hardy
temperate zone orchids from either seeds or meristem techniques on a commercial
basis. However, I suspect that after an initial burst of popular enthusiasm, any market
would decline as the home gardener encountered repeated failures.

Orchids are fertilized by insects and often the orchid requires a specific species
of insect. There is the famous example of the orchid with a ten inch spur for which
Darwin predicted a moth pollinator with a ten inch tongue. The moth was discovered
later. There are several Orphrys where the lip becomes bulbous and intricately
marked to resemble a female bee. The male alights and attempts copulation which
fertilizes, not the bee, but the flower.

Orchid seed is the smallest of all seeds with only a few hundred cells in each
seed. A capsule of Cypripedium will contain ten to twenty thousand seeds. These are
dispersed by wind which is aided by a thin membrane in which the seed is embedded.
Survival is by massive numbers of dispersed seed with a minutely few maturing.

A successful cycle in the wild starts with a fungal invasion of the seed. The
mycelia of the fungi invade the orchid seed, but the orchid seed contains chemicals
that kill the mycelia but only after they have made considerable growth. The dead
fungal mycelia provide the nutrient for the orchid seedling and possibly necessary
metabolites that the orchid cannot make for itself. This is termed symbiotic germination
but in fact seems to benefit only the orchid. The orchid cells grow and divide to
produce a protocorm. With species like Calopogon the protocorm is nearly spherical
but with others an irregular shape develops with roots ultimately emerging. With
saprophytic orchids the plant does not progress beyond this stage underground and
flowering stems emerge directly from this protocorm. Even with non-saprophytic
orchids, the first growth above ground is either a flowering stem or a mature true leaf.

It was a great breakthrough in 1922 when L. Knudson developed agar gel
compositions in which a variety of tropical orchids would germinate and grow to
mature plants. One of the key ingredients was common sugar, sucrose. Currently
propagation by meristem culture dominates, but the Knudson method was the source
of market supply for many years. The Knudson method was tried many times with
temperature zone orchids, always with failure until recently.

In the eighties a second breakthrough occurred when several investigators
reported the propagation of hardy orchids from seed using agar gel compositions of 20-
30 components coupled with introduction of certain fungi. This is termed symbiotic
germination. Germination and growth will occur without the fungi, but at present the
fungi give much superior results. Geoffrey Hadley has tested a variety of fungi and so
many were successful as to suggest that orchids are not species specific for the fungi,
although it can hardly be expected that any fungal species will work. Vivid evidence of
success was presented by W. W. Ballard, James Coke; Ian Donovan, Geoffrey Hadley,
Robert B. Mitchell, Warren Stoutamire, Carson Whitlow, and Robert 'A. Yannetti.
Yannetti exhibited a small glass jar with a screw type cover in which four Arethusa
bulbosa were growing on a half inch thick layer of gel with two of the four having.a
perfect flower and the other two in bud. This species is nearly, saprophytic as the
green stem has only a few bracts and most of the nutrition must come from non-
photosynthetic sources.Generally the seeds were disinfected with a 0.5\% NaOCI
(sodium hypOchlorite) prepared by diluting commercial bleach 1:10. The seeds float
and Coke described evacuating a vacuum chamber which removed the air from the
seeds after.which the seeds would sink in water to which surfactant had been added.
An easier technique was described by Mitchell. This was to disinfect the seed capsule
just before it splits and to handle the seed under sterile conditions.
When everything is optimum, true leaves appear in 6-18 months and flowering
was achieved in 2-6 years. Extensive losses can occur on transferring the seedlings
with leaves out of the agar.gel, but this is a problem of culture. Although the fungi are
requisite and beneficial to the germination and development of the seedlings, it is not
evident that they are necessary to mature plants except for saprophytes. Fungal
mycelia invade only the outer half of the long roots and the infections come and go.
Regular treatments with fungicide as practiced by Paul Keisling is not only beneficial
but virtually requisite suggesting that the mature plants do not require the action of
fungi. Orchid seed is long lived in dry storage. Stoutamire reported that three year old
seed of Cypripedium reginae and ten year old seed of Cypripedium acaule
germinated as well as fresh seed. There is some evidence that dry storage is required
for good germination. The following recent references deal with germination of hardy
orchids and were selected from an exhaustive bibliography that was distributed at the
Brandywine Symposium. A. B. Anderson, Wildflowers: Canada's Natl. Mag. Wild Flora
.1, 23 (1985); J. Arditti, J. Orchid Biology: Rev, and Perspectives II, Cornell Univ., New
York 1982, p.243; Am. J. Botany L4.2 442 (1981); Am. Orchid Soc. Bull. 0, 162
(1982); J. Arditti and A. Oliva, Bot. Gazette 145., 495 (1984); W. W. Ballard, Am. Orchid
Soc. Bull. fifi, 935 (1987); W. Frosch, Am. Orchid Soc. Bull. a 14 (1986); J. E.
Henrich, J. Am. Soc. Hort. Sd. 106,193 (1981).
Since the Symposium there has been continued activity in this field. A recent
paper by Ben Linden (Ann. Bot. Fennici 29:305-313, 1992) describes various
techniques for improving germination. These include treatments with disinfectants,
intermittent vacuum to introduce water into the testa, one week soaking in water in a
rotary shaker, and enzyme treatment to degrade the cell walls. This paper
summarizes the latest techniques including the components of the medium.

